\clearpage
\chapter*{Errata da exposição n\textsuperscript{o} 149}
\addcontentsline{toc}{chapter}{Errata da exposição n\textsuperscript{o} 149}
\setcounter{footnote}{0}
\markboth{Errata da exposição 149}{Errata da exposição 149}
\sid{e149}
\src{298}

\noindent\textbf{1. Teorema de dualidade para os feixes algébricos
coerentes.}

\begin{center}
[Seminário Bourbaki, t. 9, 1956/57, n\textsuperscript{o} 149, 25 p.]
\end{center}

\noindent\textbf{Observação sobre a página 14.} --- Como assinalei em minha
conferência no Congresso Internacional de Matemáticos de 1958,\footnote{
Grothendieck (Alexander). --- The cohomology theory of abstract algebraic
varieties, \emph{Proceedings of the International Congress of Mathematicians}
[1958, Edinburgh], p. 103--118. --- Cambridge, at the University Press, 1960.}
as questões levantadas aqui estão agora completamente resolvidas.

O leitor encontrará mais informações sobre a dualidade dos feixes
coerentes em \emph{loco citato}\footnotemark[1], p. 112--115, em EGA III,
2\textsuperscript{a} parte (em preparação) e, por enquanto, em SGA 1962.
Um tratamento mais sistemático encontrar-se-á em um capítulo posterior de EGA
(capítulo IX no plano previsto).
% Explicit module paragraph boundary; no prose change.
\par

\par
